% New post-checkpoint-14 full-generation criterion and toric application.
% Complete author/root/B/C actual-file and current primary source audits passed.
\begin{proposition}[Proved: a projective model with no additional divisors]
\label{D:codimension-one-projective-model}
Let $S$ be a normal proper connected scheme over an algebraically closed
field. Suppose that there is a normal projective integral scheme $P$
and an isomorphism of nonempty opens
\[
                  \phi:S^\circ\xrightarrow{\ \sim\ }P^\circ
\]
such that $P\setminus P^\circ$ has codimension at least two.
No codimension condition is imposed on $S\setminus S^\circ$.
Then, for every
proper closed subset $T\subset S$, the original-boundary tame group
$Q(S,T)$ is topologically finitely generated. In particular, if $T$
has codimension at least two, the full fundamental group
$\pi_1(S\setminus T)$ is topologically finitely generated.

More generally, D.2(ii) holds for a proper scheme whose normalized
irreducible components each admit such a projective model, for every
closed omitted set of codimension at least two in every component it
meets.
\end{proposition}

\begin{proof}
Put $U=S\setminus T$ and
\[
        V=S^\circ\cap U,\qquad V_P=\phi(V),\qquad B=P\setminus V_P.
\]
The schemes $V$ and $V_P$ are isomorphic nonempty opens, and $B$ is
closed. Every codimension-one point $\eta$ of $B$ belongs to
$P^\circ$, since $P\setminus P^\circ$ has codimension at least two.
Its corresponding point in $S^\circ$ has codimension one in $S$ and
belongs to $T$. The local discrete valuation rings at these two
points are identical under $\phi$.

A connected finite etale cover of the normal integral scheme $U$ is
normal and integral. It remains connected on the dense open $V$.
Hence restriction gives a surjection
\[
                     \pi_1(V_P)\twoheadrightarrow\pi_1(U).
\]
If this cover is classified by $Q(S,T)$, then its restriction is
tame at every codimension-one point of $B$, by the equality of
the corresponding boundary DVRs. Consequently the composite map
to $Q(S,T)$ factors through a surjection
\[
                         Q(P,B)\twoheadrightarrow Q(S,T).
                         \tag{D.CP.1}
\]
The source is finitely generated by the projective case of
Proposition~\ref{D:fixed-boundary-positive}. This proves the first
assertion. If $T$ has codimension at least two, no original
codimension-one inertia groups are killed, so $Q(S,T)$ is the full
fundamental group of $U$.

For the final assertion, pass to the reduction and then to the
disjoint union of the finite normalizations of its irreducible
components. The inverse image of the omitted set has codimension
at least two on each such component, because finite morphisms
preserve dimensions of closed irreducible subsets. The proved
normal case gives finite generation on each normalized open.
Finite proper effective descent and its finite-generation
consequence, SGA~1, Expose~IX, Theorem~4.7 and Corollary~5.2,
give finite generation downstairs
\citep[pp.~238--239 and 247--248]{SGA1}.
\end{proof}

\begin{corollary}[Proved: arbitrary boundaries on complete toric varieties]
\label{D:complete-toric-fixed-boundary}
Let $S$ be a complete toric variety over an algebraically closed field
of arbitrary characteristic. For every proper closed subset
$T\subset S$, without a torus-invariance assumption on $T$,
$Q(S,T)$ is topologically finitely generated. Thus D.1 holds for
every complete toric compactification, including nonprojective ones.
If $T$ has codimension at least two, the full group
$\pi_1(S\setminus T)$ is finitely generated.

The corresponding D.2(ii) conclusion also holds for a proper scheme
whose normalized irreducible components are complete toric varieties.
\end{corollary}

\begin{proof}
Fujino and Sato prove that every complete toric variety has a
projective $\mathbf Q$-factorial toric model which is isomorphic to it
in codimension one. Their ground-field convention is explicitly
algebraically closed of arbitrary characteristic
\citep[Theorem~1.1, p.~1; Theorem~1.3, p.~2; full proof in
Section~3, pp.~6--7]{FujinoSato2025D}.

We spell out why the result supplies the common opens required in
Proposition~\ref{D:codimension-one-projective-model}. The toric
birational map identifies the dense tori and their lattices. Being
an isomorphism in codimension one means that the two fans have the
same one-dimensional cones. The common subfan consisting of the
zero cone and these rays defines isomorphic open toric subvarieties
on the two sides. Their complements are unions of orbits of
codimension at least two. Equivalently, this is the common-open
meaning of isomorphism in codimension one used explicitly by the
authors immediately after Theorem~1.1. The model is normal because
it is a toric variety.

Apply Proposition~\ref{D:codimension-one-projective-model} with this
projective model. Its proof applies to every closed $T$; it does not
use invariance of the boundary or introduce extra divisorial
valuations. The nonnormal statement is the final assertion of the
same proposition.
\end{proof}

\begin{remark}[Scope of this full-generation conclusion]
The first proposition requires a dense common open whose complement
on the projective side has codimension at least two. Divisors may be
removed on the $S$ side: removing them does not create additional
boundary DVRs on the projective side. No morphism from $S$ to its
projective model is required. The criterion therefore includes
Proposition~\ref{D:generically-finite-projective-target} after its
Stein factorization, and also the possibly rational codimension-one
toric comparison used above.
The toric source proves this condition in all dimensions. It does
not establish it for arbitrary proper varieties: the source states
the general threefold projectivization question separately as
Problem~1.2. Accordingly the present full-generation class does not
close the unrestricted proper case of D.1 or D.2(ii).
\end{remark}
